\subsection{Mise en oeuvre du script de récupération des données}
\subsubsection{Création du compte de service}

Créer un compte de service dédié à la fonction de dépôt et le rendre propriétaire du répertoire hébergeant les dépôts.
\begin{lstlisting}[language=bash, caption=Compte de service]
useradd svc_repodnf -s /sbin/nologin
chown -R svc_repodnf /var/www/html /opt/sync_sources
\end{lstlisting}

\subsubsection{Script de récupération des sources}
Le script doit être en lecture seule, ceci afin d'éviter toute modification par un éventuel pirate ayant réussi à s'insérer dans le système.
\begin{lstlisting}[language=bash, caption=Script de récupération des sources]
echo -e '#!/bin/sh

# This script is part of Infr@Home Information System
# It is under GPL-V3 license
#
# Versionning
# YYYYMMDD_hhmm | Author                | Changelog
# 20201018_1436 | jmy37                 | Script creation
# 20201210_1431 | jmy37                 | Add rsync and wget repositories
# 20220128_1736 | jmy37                 | Add Rocky Linux repositories
# 20221107_2104 | jmy37                 | Add Alma Linux repositories
# 20240525_0959 | jmy37                 | Remove CentOS and Rocky repositories
#

# Variables
WEB_PATH=/var/www/html
LOG_PATH=/var/log
ERROR_LOG=$LOG_PATH/sync_error.log
DEBUG_LOG=$LOG_PATH/sync_debug.log
REPOLIST_FILES=/opt/sync_sources
ALMA_REPOLIST=$REPOLIST_FILES/alma
RSYNC_REPOLIST=$REPOLIST_FILES/rsync
WGET_REPOLIST=$REPOLIST_FILES/wget

# System variables
TOUCH_BIN=/usr/bin/touch
REPOSYNC_BIN=/usr/bin/reposync
RSYNC_BIN=/usr/bin/rsync
WGET_BIN=/usr/bin/wget
CURL_BIN=/usr/bin/curl

# Create log files
$TOUCH_BIN $ERROR_LOG
$TOUCH_BIN $DEBUG_LOG

# Starting synchronization
NOW=$(date +"%Y %m %d  %T")
echo -e "["$NOW"] Start syncing..." >> $DEBUG_LOG

# Alma Linux sync
while IFS=, read -r ALMA_RELEASE REPOID
do
  NOW=$(date +"%Y %m %d  %T")
  echo -e "["$NOW"] Starting to sync ["$REPOID"]..." >> $DEBUG_LOG
  if $REPOSYNC_BIN --download-metadata --delete --download-path=$WEB_PATH/alma/$ALMA_RELEASE/ --repoid=$REPOID >> $DEBUG_LOG ; then
    NOW=$(date +"%Y %m %d  %T")
    echo -e "["$NOW"] Repo ["$REPOID"] successfully synced" >> $DEBUG_LOG
  else
    echo -e "["$NOW"] Error syncing ["$REPOID"]" >> $ERROR_LOG
    echo -e "["$NOW"] Error syncing ["$REPOID"]" >> $DEBUG_LOG
  fi
done < $ALMA_REPOLIST

# Rsync repositories
while IFS=, read -r REPOID RSYNC_URL
do
  NOW=$(date +"%Y %m %d  %T")
  echo -e "["$NOW"] Starting to sync ["$REPOID"]..." >> $DEBUG_LOG
  if $RSYNC_BIN --archive --hard-links --numeric-ids --stats rsync://$RSYNC_URL $WEB_PATH/$REPOID/ >> $DEBUG_LOG ; then
    NOW=$(date +"%Y %m %d  %T")
    echo -e "["$NOW"] Repo ["$REPOID"] successfully synced" >> $DEBUG_LOG
  else
    echo -e "["$NOW"] Error syncing ["$REPOID"]" >> $ERROR_LOG
    echo -e "["$NOW"] Error syncing ["$REPOID"]" >> $DEBUG_LOG
  fi
done < $RSYNC_REPOLIST

# Wget repositories
while IFS=, read -r REPOID WGET_CUTDIRS WGET_DOMAIN WGET_URL
do
  NOW=$(date +"%Y %m %d  %T")
  echo -e "["$NOW"] Starting to sync ["$REPOID"]..." >> $DEBUG_LOG
  if $WGET_BIN --quiet --reject=md5,txt,html,tmp --recursive --no-host-directories --cut-dirs=$WGET_CUTDIRS --convert-links --no-parent --domains=$WGET_DOMAIN --directory-prefix=$WEB_PATH/$REPOID/ $WGET_URL >> $DEBUG_LOG ; then
    NOW=$(date +"%Y %m %d  %T")
    echo -e "["$NOW"] Repo ["$REPOID"] successfully synced" >> $DEBUG_LOG
  else
    echo -e "["$NOW"] Error syncing ["$REPOID"]" >> $ERROR_LOG
    echo -e "["$NOW"] Error syncing ["$REPOID"]" >> $DEBUG_LOG
  fi
done < $WGET_REPOLIST

# End of script
echo -e "["$NOW"] Syncing is over. Thanks to check log above for errors." >> $DEBUG_LOG

exit\n' > /opt/sync_sources/sync_sources.sh

chmod 500 /opt/sync_sources/sync_sources.sh
\end{lstlisting}

\subsubsection{Création des chemins indispensables}
Les chemins de base doivent être créés manuellement.
\begin{lstlisting}[language=bash, caption=Création des chemins]
mkdir -p /var/www/html/alma/8
mkdir -p /var/www/html/alma/9
mkdir -p /var/www/html/clamav-antivirus
mkdir -p /var/www/html/dataiku-dss
mkdir -p /var/www/html/geoserver
mkdir -p /var/www/html/itop
mkdir -p /var/www/html/openstreetmap/data
mkdir -p /var/www/html/openstreetmap/src
mkdir -p /var/www/html/opnsense
mkdir -p /var/www/html/redmine
mkdir -p /var/www/html/rpm-gpg-key
mkdir -p /var/www/html/yum.repos.d.sample
\end{lstlisting}

Les fichiers sources doivent également être créés manuellement.
\begin{lstlisting}[language=bash, caption=Création des fichiers sources]
touch /opt/sync_sources/alma
touch /opt/sync_sources/rsync
touch /opt/sync_sources/wget
\end{lstlisting}

Les fichiers doivent être la propriété du compte de service.
\begin{lstlisting}[language=bash, caption=Propriété des fichiers sources]
chown -R svc_repodnf: /opt/sync_sources/
\end{lstlisting}

Les fichiers de logs doivent enfin eux aussi être créés manuellement.
\begin{lstlisting}[language=bash, caption=Création des fichiers de logs]
touch /var/log/sync_debug.log /var/log/sync_error.log
chown svc_repodnf: /var/log/sync_*.log
\end{lstlisting}

La totalité du répertoire de synchronisation doit être la propriété du compte de service.
\begin{lstlisting}[language=bash, caption=Propriété du répertoire de synchronisation]
chown -R svc_repodnf: /var/www/html/
\end{lstlisting}

\subsubsection{Planification par cron}
La tâche planifiée peut s'exécuter tous les jours pour garantir la fraicheur des données.
\begin{lstlisting}[language=bash, caption=Planification par cron]
echo -e "  0  0  *  *  * svc_repodnf /opt/sync_sources/sync_sources.sh" >> /etc/crontab
\end{lstlisting}